\documentclass[11pt,a4paper]{article}
\usepackage[T1]{fontenc}
\usepackage[utf8]{inputenc}
\usepackage{lmodern,microtype,amsmath,amssymb,amsthm,mathtools}
\usepackage[left=26mm,right=26mm,top=24mm,bottom=25mm,headheight=14pt]{geometry}
\usepackage{booktabs,array,longtable,xcolor,fancyhdr}
\usepackage[hypertexnames=false]{hyperref}
\definecolor{ink}{HTML}{214F64}
\hypersetup{colorlinks=true,linkcolor=ink,citecolor=ink,urlcolor=ink,
pdftitle={Perfect Skewness of Free Stable Densities},
pdfauthor={Research manuscript prepared with ChatGPT for Keio},
pdfsubject={A complete analytic proof draft addressing Hasebe--Simon--Wang, Remark 13}}
\pagestyle{fancy}
\fancyhf{}
\lhead{\small\itshape Perfect skewness of free stable densities}
\rhead{\small 8 November 2026}
\cfoot{\thepage}
\renewcommand{\headrulewidth}{0.3pt}
\newtheorem{theorem}{Theorem}[section]
\newtheorem{lemma}[theorem]{Lemma}
\newtheorem{proposition}[theorem]{Proposition}
\newtheorem{corollary}[theorem]{Corollary}
\theoremstyle{definition}
\newtheorem{definition}[theorem]{Definition}
\theoremstyle{remark}
\newtheorem{remark}[theorem]{Remark}
\numberwithin{equation}{section}
\newcommand{\R}{\mathbb{R}}
\newcommand{\Cplus}{\mathbb{C}^{+}}
\newcommand{\sign}{\operatorname{sgn}}
\newcommand{\supp}{\operatorname{supp}}
\newcommand{\hp}{h_\ast}
\setlength{\emergencystretch}{2em}
\setlength{\parskip}{0.25em}
\allowdisplaybreaks[1]
\title{\textbf{Perfect Skewness of Free Stable Densities}\\[0.4em]
\large An equal-level curvature comparison and a global midpoint theorem}
\author{Research manuscript prepared with ChatGPT for Keio}
\date{8 November 2026\\[0.3em]
\small Version 1.0: complete analytic proof draft for external review}

\begin{document}
\maketitle

\begin{abstract}
We prove that every nondegenerate free stable density is perfectly skew:
the difference between its values at two points equidistant from its mode
has a constant sign on the entire positive half-line. For the strictly
stable normalization with parameters $(\alpha,\rho)$, the sign is exactly
$\sign((1-\alpha)(\rho-\tfrac12))$. This gives an affirmative answer to
the question in Hasebe--Simon--Wang, Remark 13 (2020), and also determines
the direction of the comparison. We obtain quantitative inequalities for
the midpoints of all density level sets. A separate logarithmic argument
treats every non-strict free stable law of index one.
The main estimate compares the curvature of an angular log-density
potential at equal levels. Integration of squared slopes transfers this
comparison to an exact differential identity for the physical midpoint.
The proof is analytic throughout, including the extremal admissible
parameters and finite support edges. The boundary coordinates and basic
unimodality ingredients are credited to prior work. Historical priority
is not inferred merely from the absence of a solution in our literature
search.
\end{abstract}

\noindent\textbf{Keywords.}
Free stable law; perfect skewness; level set; mode; curvature comparison.

\section{The original question and the results}
\label{sec:intro}

For a continuous unimodal density $f$ with unique mode $M$, perfect
skewness means that
\begin{equation}\label{eq:original}
 f(M+t)-f(M-t),\qquad t>0,
\end{equation}
has a constant sign, with the identically zero case allowed. The density
is extended by zero outside its support. We call the skewness strict
when the comparison is strict for every $t>0$.

Hasebe, Simon and Wang ask whether this property holds for the full
free strictly stable family in \cite[Remark 13, p.~321]{HSW}. The question
already appears as Remark 18 in the May 2018 preprint \cite{HSWpre}.
Their published remark records several special cases and leaves the
general case open. We retain exactly the density comparison
\eqref{eq:original}. In particular, the midpoint quantities introduced
below are proof devices, not replacements for the original target.

Write
\[
 G_\mu(z)=\int_{\R}\frac{\mu(\mathrm dx)}{z-x},\qquad
 F_\mu(z)=\frac1{G_\mu(z)},\qquad
 \phi_\mu(z)=F_\mu^{-1}(z)-z .
\]
The inverse is initially understood near infinity in the upper
half-plane. We use the standard free strictly stable normalization
\begin{equation}\label{eq:normalization}
 \phi_{\alpha,\rho}(z)=-e^{i\pi\alpha\rho}z^{1-\alpha},
 \qquad z\in\Cplus ,
\end{equation}
with the principal branch of the power and admissible parameters
\begin{equation}\label{eq:admissible}
 \mathcal A=
 \{0<\alpha\le1,\ 0\le\rho\le1\}
 \ \cup\
 \{1<\alpha\le2,\ 1-1/\alpha\le\rho\le1/\alpha\}.
\end{equation}
See \cite{BP,HK,HSW} for the classification and conventions. The pairs
$(1,0)$ and $(1,1)$ give point masses and are excluded from all density
statements.

\begin{theorem}[Global comparison in the original normalization]
\label{thm:strict}
For every nondegenerate admissible $(\alpha,\rho)$, let
$f_{\alpha,\rho}$ be the density associated with
\eqref{eq:normalization}, and let $M_{\alpha,\rho}$ be its unique mode.
Then, for every $t>0$,
\begin{equation}\label{eq:main}
 \boxed{\;
 \sign\!\left(
 f_{\alpha,\rho}(M_{\alpha,\rho}+t)
 -f_{\alpha,\rho}(M_{\alpha,\rho}-t)\right)
 =
 \sign\!\left((1-\alpha)(\rho-\tfrac12)\right).
 \;}
\end{equation}
In particular, the comparison is strict whenever
$\alpha\ne1$ and $\rho\ne\tfrac12$. If $\rho=\tfrac12$, or if
$\alpha=1$ and the law is nondegenerate, the density is symmetric about
its mode.
\end{theorem}

For each density considered here, every $0<h<\hp:=f(M)$ has exactly two
points of height $h$:
\[
 x_L(h)<M<x_R(h),\qquad
 f(x_L(h))=f(x_R(h))=h.
\]
Let
\begin{equation}\label{eq:midpoint-def}
 C(h)=\frac{x_L(h)+x_R(h)}2.
\end{equation}

\begin{theorem}[Quantitative midpoint comparison]
\label{thm:midpoint}
In the normalization \eqref{eq:normalization}, suppose
$0<\alpha<2$, $\alpha\ne1$, and $\rho>\tfrac12$ is admissible.
Then $M:=M_{\alpha,\rho}>0$, and, for every $0<h<\hp$,
\begin{align}
 1<\alpha<2:\quad&
 C(h)<M(h/\hp)^{\alpha-1}<M,
 \label{eq:midpoint-upper}\\
 0<\alpha<1:\quad&
 C(h)>M(h/\hp)^{\alpha-1}>M.
 \label{eq:midpoint-lower}
\end{align}
More precisely, $h^{1-\alpha}C(h)$ is strictly increasing if
$\alpha>1$ and strictly decreasing if $\alpha<1$.
Reflection gives the corresponding statements for $\rho<\tfrac12$.
\end{theorem}

The next result includes the free stable laws of index one which are
not strictly stable. The classification in \cite[Appendix]{BP},
also described in \cite[p.~298]{HSW}, reduces these laws, after
reflection and a real translation, to the logarithmic form below.

\begin{theorem}[The complete logarithmic index-one family]
\label{thm:log}
Let $\mu$ have Voiculescu transform
\begin{equation}\label{eq:log-normal}
 \phi_\mu(z)=c-iA-B\log z,\qquad
 A\ge0,\quad B>0,\quad c\in\R,\quad z\in\Cplus,
\end{equation}
where $\log$ is principal. Its density has a unique mode $M$ and
satisfies
\begin{equation}\label{eq:log-main}
 f(M+t)<f(M-t)\qquad\text{for every }t>0.
\end{equation}
Moreover its level-set midpoint satisfies
\begin{equation}\label{eq:log-bound}
 C(h)<M+B\log(h/\hp)<M,\qquad 0<h<\hp.
\end{equation}
The reflected laws have the opposite strict comparison. When $B=0$
and $A>0$, the density is a translated Cauchy density and the comparison
is equality.
\end{theorem}

\begin{corollary}[All free stable laws]\label{cor:all}
Every nondegenerate free stable density is perfectly skew about its
unique mode. The comparison is strict unless the density is symmetric
about the mode.
\end{corollary}

\begin{proof}[Reduction of Corollary \ref{cor:all} to the theorems]
The classification \cite[Appendix]{BP} gives, up to affine changes,
\eqref{eq:normalization} for indices other than one and the
Cauchy or logarithmic cases for index one. If
$g(x)=\lambda^{-1}f((x-d)/\lambda)$ with $\lambda>0$, its mode is
$m=d+\lambda M$ and
\[
 g(m+t)-g(m-t)
 =\lambda^{-1}\{f(M+t/\lambda)-f(M-t/\lambda)\}.
\]
Thus positive affine changes preserve the comparison, while reflection
reverses its direction. Theorems \ref{thm:strict} and \ref{thm:log}
therefore give the assertion.
\end{proof}

\section{An auxiliary equal-level comparison problem}
\label{sec:principle}

We first isolate the mechanism that will apply to both a sine function
and an affine function. The auxiliary problem is to compare the
curvatures of a strictly convex potential at its two points of equal
height, and then to compare the magnitudes of its two slopes.

\begin{lemma}[Equal-level curvature comparison]\label{lem:curvature}
Let $\beta>1$, $Q\ge0$, and let $L:(0,\pi)\to(0,\infty)$ be $C^2$.
Suppose
\begin{equation}\label{eq:L-hypotheses}
 (L')^2-LL''=Q,\qquad
 L(\theta)\ge\sqrt Q\,\sin\theta\quad(0<\theta<\pi).
\end{equation}
Set
\begin{equation}\label{eq:V-general}
 V(\theta)=\log L(\theta)-\beta\log\sin\theta .
\end{equation}
Then $V''\ge(\beta-1)\csc^2\theta>0$, $V$ tends to $+\infty$
at both ends, and it has a unique minimum at some $\theta_m$.
If $\theta_1<\theta_2$ satisfy $V(\theta_1)=V(\theta_2)$ and
$r=\sin\theta_2/\sin\theta_1$, then
\begin{equation}\label{eq:curvature-sign}
 \sign\bigl(V''(\theta_1)-V''(\theta_2)\bigr)=\sign(\log r).
\end{equation}
For $r>1$ the following quantitative lower bound holds, with $z=r^{-2}$:
\begin{equation}\label{eq:key-gap}
 V''(\theta_1)-V''(\theta_2)
 \ge \csc^2\theta_1
       \bigl[\beta(1-z)-(1-z^\beta)\bigr]>0.
\end{equation}
\end{lemma}

\begin{proof}
The first identity in \eqref{eq:L-hypotheses} gives
$(\log L)''=-Q/L^2$. Hence
\begin{equation}\label{eq:general-convexity}
 V''=\beta\csc^2\theta-\frac{Q}{L(\theta)^2}
 \ge(\beta-1)\csc^2\theta>0.
\end{equation}
If $Q>0$, the lower bound for $L$ gives
$V\ge\log\sqrt Q+(1-\beta)\log\sin\theta$, which tends to
$+\infty$ at each endpoint. If $Q=0$, then $\log L$ is affine
in $\theta$, so the same endpoint conclusion follows directly.
Strict convexity now gives the unique minimum.

At equal levels,
\begin{equation}\label{eq:equal-L}
 L(\theta_2)=r^\beta L(\theta_1).
\end{equation}
Suppose first that $r>1$, and put $z=r^{-2}\in(0,1)$. Then
\begin{align*}
 V''(\theta_1)-V''(\theta_2)
 &=
 \beta\csc^2\theta_1(1-z)
 -\frac{Q}{L(\theta_1)^2}(1-z^\beta)\\
 &\ge
 \csc^2\theta_1\{\beta(1-z)-(1-z^\beta)\}.
\end{align*}
The bracket is strictly positive because the function $z\mapsto z^\beta$
is strictly convex:
$z^\beta>1+\beta(z-1)$ for $z\ne1$.
If $r<1$, exchange indices and apply this argument with $1/r>1$,
using the bound for $L$ at $\theta_2$. This gives the reverse strict
comparison. If $r=1$, \eqref{eq:equal-L} gives equality of the two
curvatures.
\end{proof}

\begin{lemma}[From curvature to slope imbalance]\label{lem:slopes}
Under the assumptions of Lemma \ref{lem:curvature}, set $u_m=V(\theta_m)$
and $S=-V'$. For $u>u_m$, denote the inverse branches of $V$ by
$\theta_1(u)<\theta_m<\theta_2(u)$, and write $S_i(u)=S(\theta_i(u))$.
Suppose the sign
\[
 \sign\log\frac{\sin\theta_2(u)}{\sin\theta_1(u)}
 =\varepsilon\in\{-1,1\}
\]
is the same for every $u>u_m$. Then
\begin{equation}\label{eq:slope-imbalance}
 S_1>0>S_2,\qquad \sign(S_1+S_2)=\varepsilon
 \quad\text{at every level}.
\end{equation}
\end{lemma}

\begin{proof}
The signs of the slopes follow from strict convexity. Differentiating
along $V(\theta_i(u))=u$ gives
\begin{equation}\label{eq:squared-slope}
 \frac{\mathrm d}{\mathrm du}S_i(u)^2=2V''(\theta_i(u)).
\end{equation}
Both squared slopes tend to zero as $u\downarrow u_m$. Therefore
\begin{equation}\label{eq:squared-integral}
 S_1(u)^2-S_2(u)^2
 =2\int_{u_m}^{u}
   \{V''(\theta_1(w))-V''(\theta_2(w))\}\,\mathrm dw .
\end{equation}
The integrand is continuous at the lower limit after taking its
limit there, and has strict sign $\varepsilon$ at every larger level
by Lemma \ref{lem:curvature}. The integral has that sign as well.
Finally $S_1-S_2>0$, so
$S_1^2-S_2^2=(S_1-S_2)(S_1+S_2)$ proves the claim.
\end{proof}

\begin{lemma}[Level midpoints imply the pointwise comparison]
\label{lem:geometry}
Let $f$ be continuous, positive on the interior of an interval,
strictly increasing up to its mode $M$ and strictly decreasing after
it, and tending to zero at both ends of that interval.
Define $C(h)$ by \eqref{eq:midpoint-def}.
If the left support endpoint is $-\infty$ and $C(h)<M$ at every
$0<h<f(M)$, then $f(M+t)<f(M-t)$ for every $t>0$.
If the right support endpoint is $+\infty$ and $C(h)>M$ at every
such level, the reverse strict comparison holds.
\end{lemma}

\begin{proof}
Consider the first assertion. If $M+t$ lies in the interior of the
support, let $h=f(M+t)$. Then $x_R(h)=M+t$ and
$x_L(h)<M-t<M$. Strict increase on the left branch gives
$f(M-t)>h$. If $M+t$ is at or beyond a finite right endpoint,
its density is zero while $f(M-t)>0$, since the left endpoint is
$-\infty$. This includes every $t>0$. Reflection proves the second
assertion.
\end{proof}

\section{Boundary coordinates and the angular potential}
\label{sec:coordinates}

We now work with $0<\alpha<2$, $\alpha\ne1$. Put
\begin{equation}\label{eq:notation}
 p=\alpha-1,\qquad \beta=\alpha+1=p+2,\qquad
 A=\pi\alpha\rho,\qquad q=A-\frac{\pi\alpha}{2}.
\end{equation}
Reflection replaces $\rho$ by $1-\rho$; see \cite{HK,HSW}.
It suffices for the asymmetric case to assume $\rho>\tfrac12$,
so $0<q<\pi/2$.

\subsection{The standard boundary parametrization}

For $0<\theta<\pi$, define
\begin{equation}\label{eq:angular}
 v=A-p\theta,\qquad
 R=\left(\frac{\sin v}{\sin\theta}\right)^{1/\alpha},
 \qquad
 x(\theta)=R(\cos\theta-\sin\theta\cot v),
 \qquad
 y(\theta)=\frac{\sin\theta}{\pi R}.
\end{equation}
Admissibility gives $0<v<\pi$. These are the free stable boundary
coordinates of Biane \cite[Appendix A.1]{BP}, in the convention of
\cite{HK}. In particular, $x(\theta)$ traverses the support interior
and $f_{\alpha,\rho}(x(\theta))=y(\theta)$.
To check the formulas in our normalization, set
\[
 K(w)=w-e^{iA}w^{-p},\qquad w=R e^{i\theta}.
\]
The equation $\operatorname{Im}K(w)=0$ gives
$R^\alpha=\sin v/\sin\theta$, its real part gives $x$, and
Stieltjes inversion gives $y=-\operatorname{Im}(1/w)/\pi$.
The identification with the density, including the support endpoints,
is the standard boundary parametrization just cited.

Differentiation gives
\begin{equation}\label{eq:Rx-derivatives}
 \frac{R'}R=-\frac{\cot\theta+p\cot v}{\alpha},
 \qquad
 x'(\theta)=
 -\frac{R\sin\theta}{\alpha}
 \left\{(\cot\theta+p\cot v)^2+\alpha^2\right\}<0 .
\end{equation}
We will use this strict decrease to distinguish the physical left
branch from the angular left branch.

\subsection{A sine bound and strict convexity}

The elementary sine estimate
\begin{equation}\label{eq:sine-bound}
 \sin v\ge |p|\sin\theta
\end{equation}
holds on the whole admissible parameter range under consideration.
Indeed, if $p>0$, then
\[
 p(\pi-\theta)\le v\le\pi-p\theta.
\]
If $p<0$, write $s=-p\in(0,1)$; then
\[
 s\theta\le v\le\pi-s(\pi-\theta).
\]
In either interval the minimum of sine is at an endpoint, and
$\sin(su)\ge s\sin u$ for $0\le s\le1$, $0\le u\le\pi$,
by concavity of sine. This proves \eqref{eq:sine-bound}.

Set
\begin{equation}\label{eq:stable-VS}
 V(\theta)=\log\sin v-\beta\log\sin\theta
          =-\alpha\log(\pi y(\theta)),
 \qquad
 S(\theta)=-V'(\theta)=\beta\cot\theta+p\cot v .
\end{equation}
The function $L(\theta)=\sin(A-p\theta)$ satisfies
$(L')^2-LL''=p^2$. Thus Lemma \ref{lem:curvature} applies with
$Q=p^2$, and yields
\begin{equation}\label{eq:stable-Vpp}
 V''=\beta\csc^2\theta-p^2\csc^2v
 \ge\alpha\csc^2\theta>0.
\end{equation}
Consequently $y$ tends to zero at both angular endpoints, has a
unique maximum at $\theta_m$, and is strictly increasing before
$\theta_m$ and strictly decreasing after it. In combination with
$x'<0$, this proves strict unimodality in the physical coordinate.

The coordinates, the sine bound for $\alpha>1$, and the resulting
unimodality estimate also occur in corpus manuscript [122]
\cite[Section 2]{Corpus122}. They are not claimed as new here.
The additional ingredient is the comparison at two different
equal-level points in \eqref{eq:key-gap}, followed by
\eqref{eq:squared-integral} and the midpoint identity below.

\subsection{The mode and the support side needed globally}

At the central angle,
\[
 V'(\pi/2)=p\tan q .
\]
Since $V'$ is strictly increasing, $\theta_m<\pi/2$ if $p>0$ and
$\theta_m>\pi/2$ if $p<0$. At the mode, $S(\theta_m)=0$, so
\begin{equation}\label{eq:mode-positive}
 M=x(\theta_m)
 =\frac{2\alpha}{p}R(\theta_m)\cos\theta_m>0 .
\end{equation}

For later use, if $p>0$ and $\rho>\tfrac12$, then
$A>\pi\alpha/2>p\pi$. As $\theta\uparrow\pi$,
$\sin v$ has a positive limit, $R\to\infty$, and $x(\theta)\to-\infty$.
Thus the left support endpoint is $-\infty$. The right endpoint
is $+\infty$ unless $A=\pi$, when it is
$\alpha p^{1/\alpha-1}$.
If $p<0$ and $\rho>\tfrac12$, then $A>0$ and $A\le\pi\alpha<\pi$,
so as $\theta\downarrow0$ we have $R\to\infty$ and $x(\theta)\to+\infty$.
Thus the right support endpoint is $+\infty$. The left endpoint is
finite only when $\rho=1$, in which case it is
$\alpha(1-\alpha)^{1/\alpha-1}$.
At each finite endpoint $y(\theta)\to0$. These facts verify the
support hypotheses of Lemma \ref{lem:geometry}, including the
extremal parameters.

\section{Proof of the global strictly stable comparison}
\label{sec:strict-proof}

\subsection{Orienting the equal-level angles}

Fix $0<h<\hp$. There are exactly two angles
\[
 \theta_1<\theta_m<\theta_2,\qquad y(\theta_1)=y(\theta_2)=h .
\]
Because $x'<0$, they represent
$x_R(h)=x(\theta_1)$ and $x_L(h)=x(\theta_2)$, respectively.
We first prove
\begin{equation}\label{eq:r-orientation}
 \sign\log\frac{\sin\theta_2}{\sin\theta_1}=\sign p .
\end{equation}
For $0<t<\pi/2$, the formulas above give
\begin{equation}\label{eq:angular-reflection}
 \bigl(\pi y(\pi/2\pm t)\bigr)^\alpha
 =\frac{\cos^\beta t}{\cos(q\mp pt)}.
\end{equation}
Both denominators are positive by admissibility. Their difference is
\[
 \cos(q-pt)-\cos(q+pt)=2\sin q\sin(pt),
\]
which has the sign of $p$.

If $p>0$, it follows that $y(\pi/2+t)<y(\pi/2-t)$.
Here $\theta_1<\theta_m<\pi/2$. If
$\theta_1+\theta_2\ge\pi$, then
$\theta_2\ge\pi-\theta_1>\theta_m$, and strict decrease to the
right of the angular mode would give
\[
 y(\theta_2)\le y(\pi-\theta_1)<y(\theta_1),
\]
a contradiction. Hence $\theta_1+\theta_2<\pi$.
If $p<0$, the inequality in \eqref{eq:angular-reflection} reverses and
$\theta_2>\theta_m>\pi/2$. Were
$\theta_1+\theta_2\le\pi$, then
$\theta_1\le\pi-\theta_2<\theta_m$ and strict increase on the left
would give
$y(\theta_1)\le y(\pi-\theta_2)<y(\theta_2)$, again a contradiction.
Thus $\theta_1+\theta_2>\pi$ in this case. Finally,
\[
 \sin\theta_2-\sin\theta_1
 =2\cos\frac{\theta_1+\theta_2}{2}
       \sin\frac{\theta_2-\theta_1}{2}
\]
proves \eqref{eq:r-orientation}.

\subsection{A strict curvature gap at every level}

Equality of the density heights gives the exact relation
\begin{equation}\label{eq:equal-sines}
 \sin v_2=r^\beta\sin v_1,\qquad
 r=\frac{\sin\theta_2}{\sin\theta_1}.
\end{equation}
For $r>1$, Lemma \ref{lem:curvature} specializes to
\begin{equation}\label{eq:stable-gap}
 \begin{split}
 V''(\theta_1)-V''(\theta_2)
 &\ge\csc^2\theta_1
 \left[\beta(1-r^{-2})-(1-r^{-2\beta})\right]\\
 &>0 .
 \end{split}
\end{equation}
For $r<1$, exchange the indices as in that lemma. Together with
\eqref{eq:r-orientation}, this proves
\[
 \sign\bigl(V''(\theta_1)-V''(\theta_2)\bigr)=\sign p
\]
at every nonmaximal density level. Lemma \ref{lem:slopes} now gives
\begin{equation}\label{eq:stable-slope-sign}
 S_1>0>S_2,\qquad \sign(S_1+S_2)=\sign p,
 \qquad S_i=S(\theta_i).
\end{equation}
This is the strict global estimate needed to control the physical
midpoints. It does not require an expansion near the mode or in a tail.

\subsection{The exact midpoint differential identity}

From \eqref{eq:stable-VS},
$y'=yS/\alpha$. Combining this identity with
\eqref{eq:Rx-derivatives}, along either branch of height $h$ we get
\begin{equation}\label{eq:dx-dh}
 \frac{\mathrm dx}{\mathrm dh}
 =-\frac{D(\theta)}{\pi h^2},\qquad
 D(\theta)=
 \frac{\sin^2\theta\{(\cot\theta+p\cot v)^2+\alpha^2\}}{S(\theta)}.
\end{equation}
Since $\cot\theta+p\cot v=S-\alpha\cot\theta$, division gives
\begin{align}
 D(\theta)
 &=\sin^2\theta\{S-2\alpha\cot\theta\}
     +\frac{\alpha^2}{S}\notag\\
 &=p\sin^2\theta(\cot v-\cot\theta)+\frac{\alpha^2}{S}
  =-p\pi h\,x(\theta)+\frac{\alpha^2}{S}.
 \label{eq:D-identity}
\end{align}
Here the last equality uses
$\pi h\,x(\theta)=\sin^2\theta(\cot\theta-\cot v)$.
Averaging \eqref{eq:dx-dh} over the two branches yields the exact
identity
\begin{equation}\label{eq:midpoint-ode}
 \boxed{\quad
 C'(h)-\frac{p}{h}C(h)
 =-\frac{\alpha^2}{2\pi h^2}
       \left(\frac1{S_1}+\frac1{S_2}\right).
 \quad}
\end{equation}
Because $S_1S_2<0$, equation \eqref{eq:stable-slope-sign} shows that
the right-hand side has strict sign $p$. Multiplication by the positive
integrating factor $h^{-p}$ gives
\begin{equation}\label{eq:weighted-monotonicity}
 \sign\frac{\mathrm d}{\mathrm dh}\{h^{-p}C(h)\}=\sign p.
\end{equation}

As $h\uparrow\hp$, both inverse branches converge to the mode and
$C(h)\to M$. Strict monotonicity in \eqref{eq:weighted-monotonicity}
and this finite endpoint limit imply
\[
 \begin{cases}
 C(h)<M(h/\hp)^p,&p>0,\\
 C(h)>M(h/\hp)^p,&p<0.
 \end{cases}
\]
The inequalities are strict: between any fixed $h<\hp$ and the
endpoint there is another interior level at which strict monotonicity
already gives a positive difference.
Combining these inequalities with $M>0$ from
\eqref{eq:mode-positive} proves Theorem \ref{thm:midpoint}.
No integration across a possibly singular branch derivative at
$h=\hp$ is needed.

\subsection{Every distance, all parameters, and the degenerate corners}

For $\rho>\tfrac12$ and $\alpha>1$, \eqref{eq:midpoint-upper} gives
$C(h)<M$ at every level, and the left support endpoint is $-\infty$.
Lemma \ref{lem:geometry} gives
$f(M+t)<f(M-t)$ for every $t>0$, including points beyond a finite
right endpoint. For $\rho>\tfrac12$ and $\alpha<1$,
\eqref{eq:midpoint-lower} and the infinite right endpoint give
$f(M+t)>f(M-t)$ for every $t>0$. Reflection proves the cases
$\rho<\tfrac12$.

When $\rho=\tfrac12$, the law is symmetric. At $\alpha=1$ and
$0<\rho<1$, the transform \eqref{eq:normalization} gives
\[
 f_{1,\rho}(x)=
 \frac{\sin(\pi\rho)}
 {\pi\{(x+\cos(\pi\rho))^2+\sin^2(\pi\rho)\}},
\]
which is symmetric about its mode $-\cos(\pi\rho)$.
The two remaining index-one corners are point masses, already
excluded. At $\alpha=2$ the only admissible value is $\rho=\tfrac12$,
and the law is the symmetric semicircle. This proves
Theorem \ref{thm:strict} on its entire stated domain.
\hfill$\square$

\section{The logarithmic index-one family}
\label{sec:log-proof}

\subsection{Coordinates and support}

For \eqref{eq:log-normal}, put
\begin{equation}\label{eq:log-coordinates}
 \begin{gathered}
 T(\theta)=A+B\theta,\qquad
 R(\theta)=\frac{T(\theta)}{\sin\theta},\\
 x(\theta)=T(\theta)\cot\theta-B\log R(\theta)+c,\qquad
 y(\theta)=\frac{\sin^2\theta}{\pi T(\theta)},\qquad 0<\theta<\pi.
 \end{gathered}
\end{equation}
We use the boundary method of \cite[Appendix, Proposition A1.3]{BP},
and fix the logarithmic sign directly from the transform in
\eqref{eq:log-normal}. Indeed, with $w=Re^{i\theta}$ and
$K(w)=w+\phi_\mu(w)$,
\[
 \operatorname{Im}K(w)=R\sin\theta-A-B\theta,\qquad
 \operatorname{Re}K(w)=R\cos\theta-B\log R+c .
\]
Setting the imaginary part to zero, taking the real part, and applying
Stieltjes inversion gives exactly \eqref{eq:log-coordinates}.
Direct differentiation then gives
\begin{equation}\label{eq:log-xprime}
 x'(\theta)=
 -T(\theta)\left\{\left(\cot\theta-\frac{B}{T(\theta)}\right)^2+1\right\}
 <0.
\end{equation}
The density identification can also be justified directly. Define
\[
 \Omega=\{Re^{i\theta}:0<\theta<\pi,\ 
                   R>(A+B\theta)/\sin\theta\}.
\]
This is connected, $K(\Omega)\subset\Cplus$, and
$|w|>B$ throughout $\Omega$. Every finite boundary point maps to the
real line, and $K(w)=w+O(\log|w|)$ at infinity. Hence
$K:\Omega\to\Cplus$ is proper: a sequence leaving all compact
subsets can neither approach a finite boundary point nor tend to
infinity while its image stays in a compact subset of $\Cplus$.
Its image is open and closed in $\Cplus$, hence is all of $\Cplus$.
Since $K'(w)=1-B/w$ has no zero in $\Omega$, this proper map is an
unbranched covering of the simply connected upper half-plane, and is
therefore a conformal bijection. Its inverse agrees near infinity
with $F_\mu$ and hence agrees throughout $\Cplus$.
The boundary formulas and Stieltjes inversion thus identify the
density without importing a sign convention for the logarithmic
term from another displayed parametrization.

The left endpoint is always $-\infty$. The right endpoint is
$+\infty$ if $A>0$, and
\begin{equation}\label{eq:log-edge}
 B-B\log B+c
\end{equation}
if $A=0$. In the latter case the remaining real boundary ray $w\ge B$
maps to $[B-B\log B+c,\infty)$ under
$w\mapsto w-B\log w+c$ and has zero density.
The curved boundary in \eqref{eq:log-coordinates} covers the support
interior. In both cases $y$ tends to zero at each end.

\subsection{The same curvature comparison}

Here take $L(\theta)=T(\theta)$, $Q=B^2$, and $\beta=2$ in
Lemma \ref{lem:curvature}. Its hypotheses hold because
\[
 (T')^2-TT''=B^2,\qquad
 T(\theta)=A+B\theta\ge B\theta\ge B\sin\theta .
\]
Thus
\begin{equation}\label{eq:log-VS}
 V=\log T-2\log\sin\theta=-\log(\pi y),\qquad
 S=-V'=2\cot\theta-\frac{B}{T},\qquad
 V''=2\csc^2\theta-\frac{B^2}{T^2}>0.
\end{equation}
There is a unique angular maximum of $y$, and
\eqref{eq:log-xprime} gives a unique physical mode $M$.

At equal density height, with
$\theta_1<\theta_m<\theta_2$,
\[
 \frac{T(\theta_2)}{T(\theta_1)}
 =\left(\frac{\sin\theta_2}{\sin\theta_1}\right)^2=r^2.
\]
Since $B>0$ and $\theta_2>\theta_1$, the left-hand side exceeds one,
so $r>1$ at every level. The curvature gap has the particularly
simple form
\begin{equation}\label{eq:log-gap}
 V''(\theta_1)-V''(\theta_2)
 \ge\csc^2\theta_1(1-r^{-2})^2>0.
\end{equation}
Lemma \ref{lem:slopes} gives $S_1>0>S_2$ and $S_1+S_2>0$.

\subsection{A logarithmic midpoint bound and reflection}

Since $y'=yS$, equations \eqref{eq:log-coordinates} and
\eqref{eq:log-xprime} give
\[
 \frac{\mathrm dx}{\mathrm dh}
 =-\frac{D(\theta)}{\pi h^2},\qquad
 D(\theta)=
 \frac{\sin^2\theta\{(\cot\theta-B/T)^2+1\}}{S}.
\]
Now $\cot\theta-B/T=S-\cot\theta$, so division gives
\begin{equation}\label{eq:log-D}
 D(\theta)
 =\sin^2\theta(S-2\cot\theta)+\frac1S
 =-B\pi h+\frac1S .
\end{equation}
Consequently
\begin{equation}\label{eq:log-ode}
 C'(h)
 =\frac Bh-\frac1{2\pi h^2}
          \left(\frac1{S_1}+\frac1{S_2}\right)
 >\frac Bh .
\end{equation}
It follows that $C(h)-B\log h$ is strictly increasing. Taking its
finite limit $M-B\log\hp$ at the top level gives
$C(h)<M+B\log(h/\hp)<M$. The infinite left endpoint and
Lemma \ref{lem:geometry} now give the strict comparison for every
$t>0$, including the finite right edge when $A=0$.

For clarity, reflection must respect the logarithm branch. If
$\phi_X(z)=c-iA-B\log z$ with $A\ge0$, $B>0$, then
\begin{equation}\label{eq:log-reflection}
 \phi_{-X}(z)
 =-\overline{\phi_X(-\overline z)}
 =-c-i(A+B\pi)+B\log z .
\end{equation}
This follows from
$\log(-\overline z)=\overline{\log z}+i\pi$ for $z\in\Cplus$.
Every admissible logarithmic law with the opposite coefficient is
a reflection of this form. For $B=0$, $A>0$, the inverse transform
gives a translated Cauchy density. This completes the proof of
Theorem \ref{thm:log}, and hence of Corollary \ref{cor:all}.
\hfill$\square$

\section{What the argument adds, and the limits of the claim}
\label{sec:status}

\subsection{The substantive estimate}

The method uses three steps: equal-level curvature comparison,
integration of squared slopes, and an integrating factor for the
physical midpoint. The strict remainder
\[
 \beta(1-z)-(1-z^\beta)>0,\qquad 0<z<1,\quad\beta>1,
\]
is the sign-producing estimate. The angular change of variables by
itself does not yield this comparison. Nor does strict unimodality
alone determine which side of the mode has the larger density at
every distance. Theorems \ref{thm:midpoint} and \ref{thm:log}
also quantify the displacement of all equal-height interval
midpoints; no optimality of these quantitative bounds is asserted.

For index one, the usual representation
$a\mathcal C+b\mathcal T$, with $a\ge0$, $b\in\R$, where $\mathcal C$ is standard Cauchy,
$\phi_{\mathcal T}(z)=-\log z$, and the variables are freely
independent, fits \eqref{eq:log-normal}: for $b>0$ take
$A=a$, $B=b$, and $c=b\log b$, with any extra location added to $c$.
The direction of strict perfect skewness is therefore
$-\sign(b)$ when $b\ne0$. This covers arbitrary Cauchy perturbations
of the exceptional law.

\subsection{Prior work and the bounded literature audit}

Biane's density coordinates and unimodality theorem are established
background \cite{BP}. Hasebe--Kuznetsov \cite{HK} provide the transform
normalization used here. Corpus [122] \cite{Corpus122} already uses
the angular sine bound to study the mean, median and mode, so those
ingredients are explicitly credited above.
The exact target of the present proof is the general question in
\cite[Remark 13]{HSW}, not the mean--median--mode question in its
following remark.

The literature audit was carried out on 8 November 2026. It included
the primary sources \cite{BP,HK,HSW}, the May 2018 version
\cite{HSWpre}, the related 2026 preprint \cite{Wang2026}, and the
130 TEX manuscripts in the supplied corpus at the commit recorded
in \cite{Corpus122}. In the inspected material we did not locate an
earlier general perfect-skewness theorem or the same combined proof.
Wang's 2026 preprint concerns classical divisibility and density-shape
questions; its main results and relevant sections were inspected
for overlap, without using its theorems as premises of our proof.

This is a bounded audit, not a proof that no antecedent exists.
Accordingly, mathematical completeness of the argument presented
here, independent external verification, and historical priority
remain distinct claims. The manuscript is a complete proof draft
for external review; it does not report acceptance or verification
by a journal or by the original question's authors.

\subsection{The classical problem is a separate target}

The classical stable analogue is associated with Hall's paper
\cite{Hall} in the discussion of \cite[Remark 13]{HSW}.
The present proof uses the inverse \emph{free} Cauchy transform
\eqref{eq:normalization}. It provides no proof of the corresponding
classical stable-density statement. In particular, an affirmative
answer to the free question must not be described as a resolution
of Hall's classical problem.

\clearpage
\appendix
\section{Checks against the requested success conditions}
\label{sec:gates}

The following audit records what is established by the manuscript
and what the bibliographic search cannot certify.

\begingroup
\small
\renewcommand{\arraystretch}{1.15}
\begin{longtable}{@{}p{0.07\textwidth}>{\raggedright\arraybackslash}p{0.26\textwidth}>{\raggedright\arraybackslash}p{0.60\textwidth}@{}}
\toprule
Gate & Requirement & Evidence and status\\
\midrule
\endhead
G1 & A pre-existing identified problem &
The target is the explicit perfect-skewness question in
Hasebe--Simon--Wang, Remark 13 (2020), already in the 2018 preprint
as Remark 18. The reference and original density comparison are
given in Section \ref{sec:intro}.\\
G2 & The original statement is unchanged &
Theorem \ref{thm:strict} proves the same comparison for all admissible
parameters and all $t>0$, and specifies its sign. Lemma
\ref{lem:geometry} proves the implication from the auxiliary
midpoints to that original statement.\\
G3 & The problem was actually unresolved &
The original source poses the general case as open. No later general
solution was located in the inspected primary material. The bounded
audit does not certify the absence of every intervening result.\\
G4 & An unconditional complete proof &
Sections \ref{sec:principle}--\ref{sec:log-proof} close the analytic
argument for all stated parameters and levels, including the support
edges. No conjecture, numerical experiment, or finite-height
exclusion is a premise. External mathematical review remains pending.\\
G5 & A substantive step beyond reparametrization &
The strict equal-level curvature gap \eqref{eq:key-gap}, the
squared-slope integral \eqref{eq:squared-integral}, and the midpoint
identities \eqref{eq:midpoint-ode} and \eqref{eq:log-ode} supply the
global sign. The earlier coordinates and sine bound are credited.\\
G6 & No earlier identical solution &
No duplicate theorem with the same general scope or essentially
identical argument was found in the audited sources. Historical
priority is not certified by search absence.\\
\bottomrule
\end{longtable}
\endgroup

\noindent\textbf{Quantifier audit.}
The proof includes $0<\alpha<1$ and $1<\alpha<2$, both extremal
admissible skewness parameters, the symmetric case, the nondegenerate
strictly stable index-one cases, $\alpha=2$, and every logarithmic
index-one law after reflection. All density heights $0<h<\hp$ and
all distances $t>0$ are covered. Point masses are excluded because
they have no Lebesgue density.

\noindent\textbf{Reproducibility.}
This single TEX file contains the complete proof and bibliography.
It requires no external data, figure, bibliography file, computer
algebra output, or numerical certificate. It compiles with two runs
of \texttt{pdflatex}.

\clearpage
\begin{thebibliography}{9}

\bibitem{BP}
H.~Bercovici and V.~Pata, with an appendix by P.~Biane,
\emph{Stable laws and domains of attraction in free probability theory},
Annals of Mathematics (2) \textbf{149} (1999), 1023--1060.
\href{https://doi.org/10.2307/121080}{doi:10.2307/121080}.
\href{https://arxiv.org/abs/math/9905206}{arXiv:math/9905206}.
The density parametrizations used here are in Appendix A.1.

\bibitem{HK}
T.~Hasebe and A.~Kuznetsov,
\emph{On free stable distributions},
Electronic Communications in Probability \textbf{19} (2014),
paper no.~56, 1--12.
\href{https://doi.org/10.1214/ECP.v19-3443}{doi:10.1214/ECP.v19-3443}.

\bibitem{HSW}
T.~Hasebe, T.~Simon and M.~Wang,
\emph{Some properties of the free stable distributions},
Annales de l'Institut Henri Poincar\'e, Probabilit\'es et Statistiques
\textbf{56} (2020), no.~1, 296--325.
\href{https://doi.org/10.1214/19-AIHP962}{doi:10.1214/19-AIHP962}.
\href{https://www.math.sci.hokudai.ac.jp/~thasebe/HasebeSimonWang-AIHP2020.pdf}{Author-hosted published PDF}.
The target is Remark 13, p.~321; the index-one classification is
discussed on p.~298.

\bibitem{HSWpre}
T.~Hasebe, T.~Simon and M.~Wang,
\emph{Some properties of the free stable distributions},
\href{https://arxiv.org/abs/1805.01133v2}{arXiv:1805.01133v2},
7 May 2018; first submitted 3 May 2018.
The corresponding question is Remark 18, p.~39, in this version.

\bibitem{Hall}
P.~Hall,
\emph{On unimodality and rates of convergence for stable laws},
Journal of the London Mathematical Society (2) \textbf{30} (1984),
no.~2, 371--384.
\href{https://doi.org/10.1112/jlms/s2-30.2.371}{doi:10.1112/jlms/s2-30.2.371}.

\bibitem{Wang2026}
M.~Wang,
\emph{Classical infinite divisibility, self-decomposability and
bell-shape of free stable laws},
preprint, 9 October 2026, Hexagon 2610.00169v1, 77 pp.
\href{https://hexagonmath.org/pdf/2610.00169v1}{Primary PDF}.
Inspected for overlap; no result from this preprint is assumed
in the proofs above.

\bibitem{Corpus122}
\emph{Mean, median and mode of stable laws, and a sharp gamma-mixture
criterion}, research manuscript prepared with ChatGPT for Keio,
3 November 2026, item [122] in
\emph{Verifiable Mathematical Reasoning Corpus},
repository maintained by QEEDDD.
\href{https://github.com/QEEDDD/Verifiable-Mathematical-Reasoning-Corpus/tree/c30d40119134a59eda6fa5cc6331f27c8e67d034}{Corpus snapshot: commit c30d401}.
File: \texttt{Mean\_Median\_Mode\_and\_Gamma\_Mixtures\_of\_Stable\_Laws.tex}
(with item prefix [122]).

\end{thebibliography}
\end{document}
